\section*{Paper \ref{paper:dafx}}

\textbf{\citetitle{miticObjectiveComparisonPanning2025}}

\subsubsection*{Authors}
\citefullauthors{miticObjectiveComparisonPanning2025}

\subsubsection*{Status}
Published in \citepublication{miticObjectiveComparisonPanning2025}, September 2025

\subsubsection*{Research methods}
Several \gls{dsp} techniques for determining audio spatial features were performed on a small dataset that had predictable musical characteristics. The different techniques were compared using a new objective metric.

\subsubsection*{Results and conclusions}
The paper revealed that although there are several \gls{dsp} techniques specifically designed for the extraction of recording and panning features, they are all limited in their own ways. One particular technique (\gls{spf} from \cite{tzanetakisStereoPanningFeatures2007}) was particularly promising as a candidate for use later in the doctoral phase of studies, and the others could be used with some modifications. Overall it highlighted the importance of choosing both a \gls{sip} and a distance metric that are appropriate to the overarching task, thereby addressing \rqb{}.

\subsubsection*{My contribution}
Main author, responsible for all data collection, experimental design, and software implementation. Presented the paper at the 28th International Conference on Digital Audio Effects in Ancona, Italy.